\section{The leakage certificate and the weight-only bound: proof and details}
\label{sm:leakage}

\paragraph{The three hypotheses of Theorem~\ref{thm:leakage}.}
\begin{description}
\item[(H1) Common origin.] $A$ and $A_P$ are referred to the same $E_0$. A shift common to both cancels
exactly (verified to $1.000000$ for shifts of $10^{-3}$, $10^{-2}$, $10^{-1}\,t$). If instead $A_P$ is
referred to $E_0+\delta$, then
$\lVert A-A_P^{(E_0+\delta)}\rVert_1\le\lVert A-A_P\rVert_1
+\lVert\phi_P\rVert^{2}(4/\pi)\arctan(|\delta|/2\eta)$, an additive term verified to $2\times10^{-5}$
in all $48$ cases tested and \emph{which the certificate does not control}. A ground state obtained by
Rayleigh--Ritz on $60\%$ of the $N$-electron sector at $L=8$ gives $\delta=6.7\times10^{-4}\,t$, which
adds $0.0024$ to a bound of $1.2036$ ($+0.20\%$) while multiplying the \emph{measured} error by $16$.
No cheap a posteriori control of $\delta$ is available: residual (Weinstein) and Temple
bounds~\cite{weinstein1934,temple1928} computed on
ground-state subspaces of $20$--$80\%$ of the $N$ sector at $L=6$ give $\arctan$ terms of $1.76$,
$1.64$, $1.40$ and $0.92$, all vacuous against the ceiling of $2$.
\item[(H2) Normalization and support of the $L_1$ norm.] The bound is on
$\int_{\mathbb R}|A-A_P|\,d\omega$ normalized by $\lVert\phi\rVert^{2}=\int_{\mathbb R}A\,d\omega$. A
relative error normalized by $\int_W A$ over a finite window and evaluated by quadrature --- the
convention of this work's own resource scan --- is a different quantity. Measured here,
$\lVert\phi\rVert^{2}/\int_W A=1.0100$, $1.0309$, $1.0373$ at $\eta=0.05$, $0.15$, $0.18\,t$ ($L=6$)
and $1.0110$, $1.0278$, $1.0335$ ($L=8$), with a numerator ratio of $1.0040$--$1.0115$ across the $204$
configurations of Sec.~\ref{sec:frontier}. The two conventions are never interchanged silently in this
paper.
\item[(H3) Exact reference.] $A$ is the exact spectral function of the full sector. If the reference is
itself a Lanczos reconstruction of finite depth $n_\ell$, the theorem does \emph{not} bound
$\lVert A^{(n_\ell)}-A_P^{(n_\ell)}\rVert_1$. Every measured error in this work is therefore labelled
with the depth of its reference and with that reference's own drift, and we quote no error lying within
a factor $20$ of that drift.
\end{description}

\paragraph{Proof of Theorem~\ref{thm:leakage}.}
\begin{proof}
Write $u=R(z)\phi$ and $u_P=R_P(z)\phi_P$, extended by zero outside $\operatorname{ran}P$.

\emph{(1) The split.} Since $u_P\in\operatorname{ran}P$ and $(z-H_P)u_P=\phi_P$ there,
$(z-H)u_P=zu_P-PHu_P-QHu_P=\phi_P-QHu_P$; subtracting this from $(z-H)u=\phi$ gives
\begin{equation}
u-u_P\;=\;R(z)\bigl[\,\phi_Q\;+\;QH\,R_P(z)\phi_P\,\bigr].
\label{eq:split}
\end{equation}
This is the Feshbach--L\"owdin partition of the resolvent~\cite{feshbach1958,lowdin1962}, written as an
exact identity rather than as an expansion: the first source term is the part of the probe the subspace
never held, the second is the part of the Hamiltonian that carries amplitude out of it. Only $P^{2}=P$,
$P^{\dagger}=P$, $H=H^{\dagger}$ and $\eta>0$ are used; no property of the determinant basis, of the
ranking, or of the sampler enters anywhere below.

\emph{(2) Pointwise.}
$A-A_P=(\eta/\pi)(\lVert u\rVert+\lVert u_P\rVert)(\lVert u\rVert-\lVert u_P\rVert)$ and
$\bigl|\lVert u\rVert-\lVert u_P\rVert\bigr|\le\lVert u-u_P\rVert$, so
$|A-A_P|\le(\eta/\pi)(\lVert u\rVert+\lVert u_P\rVert)\,\lVert u-u_P\rVert$ at every $\omega$.

\emph{(3) Integration.} By Cauchy--Schwarz in $\omega$,
\begin{widetext}
\begin{equation*}
\lVert A-A_P\rVert_1\;\le\;\frac{\eta}{\pi}
\Bigl[\int\bigl(\lVert u\rVert+\lVert u_P\rVert\bigr)^{2}d\omega\Bigr]^{1/2}
\Bigl[\int\lVert u-u_P\rVert^{2}d\omega\Bigr]^{1/2}.
\end{equation*}
\end{widetext}
The first bracket is at most $(\pi/\eta)^{1/2}(\lVert\phi\rVert+\lVert\phi_P\rVert)$, by Minkowski in
$L^{2}(d\omega)$ and the two sum rules. For the second, apply Minkowski to Eq.~\eqref{eq:split}: its
first term obeys the sum rule with seed $\phi_Q$, contributing exactly
$(\pi/\eta)^{1/2}\lVert\phi_Q\rVert$; its second obeys
$\lVert R(z)v\rVert\le\eta^{-1}\lVert v\rVert$, since $H$ is self-adjoint and $\operatorname{Im}z=\eta$,
and the definition~\eqref{eq:lambda} then contributes at most $(\pi/\eta)^{1/2}\Lambda_P(\eta)/\eta$.
Multiplying the three factors gives the leakage branch of~\eqref{eq:twobudget}.

\emph{(4) The other two branches.} $A$ and $A_P$ are nonnegative with known integrals, so
$\lVert A-A_P\rVert_1\le\int A+\int A_P=\lVert\phi\rVert^{2}(1+w_P)$ under no hypothesis at all; and
$\bigl|\int A-\int A_P\bigr|\le\lVert A-A_P\rVert_1$ gives the lower bound
$\lVert\phi\rVert^{2}(1-w_P)$.

\emph{(5) The closed form.} $R_P(z)\phi_P=\sum_m c_m(z-\tilde E_m)^{-1}\ket{\tilde m}$, so
\begin{equation*}
\bigl\lVert QHR_P(z)\phi_P\bigr\rVert^{2}\;=\;\sum_{m,m'}
\frac{c_m\bar c_{m'}\,\braket{g_{m'},g_m}}{(z-\tilde E_m)(\bar z-\tilde E_{m'})}\,.
\end{equation*}
With $x=\omega+E_0$ the remaining integral is
$\int dx\,[(x-\tilde E_m+i\eta)(x-\tilde E_{m'}-i\eta)]^{-1}$, which has one pole in each half
plane; closing above picks up $2\pi i/[(\tilde E_{m'}-\tilde E_m)+2i\eta]$, and multiplying by
$\eta/\pi$ gives Eq.~\eqref{eq:lambda}.
\end{proof}

\paragraph{What object is certified.} If the resolvent inside $\operatorname{ran}P$ is itself evaluated
by a Krylov recursion of depth $n_\ell$, then $\Lambda$ computed on that Krylov space bounds the error
of the depth-$n_\ell$ reconstruction \emph{and of nothing else}. The convenient inequality
$\Lambda_{\mathcal K}\ge\Lambda_S$ that would transfer the statement to the dense $A_S$ is false:
$1169$ of $6000$ random triples violate it, at a minimum ratio of $0.5258$. Every table of
Sec.~\ref{sec:frontier} carries the depth of the object it certifies for that reason. Full
reorthogonalization is mandatory throughout: without it $P_{\mathcal K}$ is not a projector --- the
orthogonality defect reaches $0.238$ at $L=8$, where the captured weight then reads $1.4868$, an
impossible value --- and $\Lambda$ is inflated thirtyfold.

\paragraph{Relation to the classical error theory, and the scope of the novelty claim.} Nothing in the
proof is new as analysis, and we say so before a referee does. The split~\eqref{eq:split} is L\"owdin
partitioning~\cite{lowdin1962}, the same algebra as the Feshbach projection of scattering
theory~\cite{feshbach1958}; and sixty years of Rayleigh--Ritz and Krylov error analysis for resolvents
and matrix functions~\cite{saad2011,golub2010,epperly2022} contains every rung of steps (2) and (3),
along with far sharper residual, Temple and Kato--Temple estimates for individual
eigenpairs~\cite{temple1928,weinstein1934,kato1949}, and the
loss-of-orthogonality analysis of Paige~\cite{paige1976} that makes step (5) conditional on
reorthogonalization, and the polynomial-restart convergence theory of Beattie, Embree and Sorensen~\cite{beattie2005}. Step (3) in particular is the reduced-basis bound on a Galerkin projection onto an arbitrary subspace---the dual norm of the residual over an inf--sup constant~\cite{rozza2008,hesthaven2016}---with that constant equal to $\eta$.

Nor is ``a posteriori'' ours to claim, and we name the prior art rather than wait to be shown it.
Stopping criteria for rational matrix functions of Hermitian matrices, computed from the residual
alone, are Frommer and Simoncini's~\cite{frommersimoncini2008}; Jia and Lv's title is this section's
claim word for word---\emph{a posteriori error estimates of Krylov subspace approximations to matrix
functions}~\cite{jialv2015}---and Frommer and Schweitzer carry that line to Stieltjes
functions~\cite{frommerschweitzer2016}; Lanczos approximations to matrix functions carry a posteriori bounds from the same Cauchy-integral argument as step (5)~\cite{chen2022lanczos,xuchen2025block}, and stochastic Lanczos quadrature carries them on an approximated spectral measure itself, in Wasserstein and Kolmogorov--Smirnov distance~\cite{chen2021slq}---a global density of states, not the seeded local measure bounded here; and five months before this work, on \emph{quantum} Krylov,
Oliveira, Ziems and Glaser introduce ``two new metrics, the imaginary and unitary filters, that
successfully assess the reliability of the obtained solutions without any knowledge of the true
eigenspectrum''~\cite{oliveira2026}. That is the framing of this section, established earlier and by others, and the priority is theirs. On the selection side, that a subspace chosen for the energy misses configurations a response needs, and a selection criterion built for the response instead, are Reinholdt, Kjellgren and Kongsted's, for classical selected configuration interaction~\cite{reinholdt2026lr}.

What we did not find in that literature is the conjunction of three things, and it is the whole of the
novelty claimed here. First, a closed form for the frequency-integrated residual $\Lambda_P(\eta)$, assembled from the Ritz pairs, on a subspace of \emph{coordinates} carrying no Krylov structure, no invariance and no spectral-gap assumption---so that the estimate
attaches to the set of determinants a sampler returned, which is not the span of the iterates of any
recursion, and Eq.~\eqref{eq:split} is evaluated on that set directly. Second, that the bound is
two-sided on the same object: the lower branch is fixed by the captured weight and the upper branch adds what the weight does not control, and Sec.~\ref{sec:frontier} uses the distance between them rather than either alone.
Third, a \emph{measured} calibration of where that bound crosses the trivial one, over sixteen
$(L,\eta)$ cells whose true errors span five decades (Sec.~\ref{sec:frontier:calib}), in place of an
asymptotic statement about the crossing. Whether such a bound is \emph{useful} at the
fractions sampling reaches is a separate question; it is not settled by the proof, and
Sec.~\ref{sec:frontier} answers it --- in the negative --- by measurement.

\paragraph{The pooled sweep of Fig.~\ref{fig:thm1iii-violation}, and its limits.}
The pooled sweep behind Fig.~\ref{fig:thm1iii-violation} contributes $621$ subspaces from five Hamiltonian families and eight construction rules, written by two independent implementations; $15$
reproduce their target to double-precision round-off and define no ratio. On the $606$ that remain the
weight-only bound fails on $468$, and restricted to the regime this work operates in, $w_S\ge0.99$, it fails
on $363$ of $363$, the mildest by a factor $2.10$ (Sec.~\ref{app:counterex}). The compliant cases are
subspaces so poor that the relative error is already of order unity, where the bound asserts nothing.
In the configuration of Fig.~\ref{fig:akwsampled} itself --- $L=8$, $\eta=0.18t$, the $3920$-dimensional
$(N{+}1)$ sector, momentum seeds $\hat c^{\dagger}_{k\uparrow}\ket{\Psi_0}$, both branches, eight
momenta, subspaces grown in order of accumulated Born weight --- the median violation factor at sector
fractions $0.50$, $0.70$, $0.85$, $0.95$ is $18.5$, $23.6$, $26.5$, $29.9$, the $64$ rows spanning
$15.1$ to $57.3$ overall and $15.1$ to $33.7$ at the fraction $0.85$ of Fig.~\ref{fig:akwsampled}. Those medians
\emph{rise} monotonically, by a factor $1.6$, while the measured error over the same range falls by
three orders of magnitude, from the $10^{-1}$ level at the lowest fraction to $1.0\times10^{-4}$ at
$0.95$: the weight-only bound fails worst exactly where the reconstruction is best, which is the opposite
of what an error bound is for. That is a defect of the functional form and not a pathology at an edge of
parameter space, and it is therefore measurable: Fig.~\ref{fig:thm1iii-violation} plots it against
sector fraction and against $\eta$, beside the replacement bound on the same subspaces.

Three properties of that sweep are declared with it, because each of them limits how it may be read.
Its ranking uses $K=16$ slices, one protocol step short of the $K=18$ of the resource scan of
Sec.~\ref{sec:prices}. Its momentum seeds have support on $3848$ to $3920$ of the $3920$ determinants and captured weight below one at every fraction, so the support floor derived below for the local seed does not apply to it; the members of a translation orbit carry equal Born weight, so the cut can fall inside a tied orbit. How far a tie-breaking can move an error was measured on a local-seed subspace at $L=6$, fraction $0.85$, whose cut fell among zero-amplitude determinants: ten tie-breakings moved the measured error by ${\pm}20\%$, which is why the violation factors are quoted to three figures at most and why the rise
across fractions is read as a direction and not as a rate. (That $\pm20\%$ is the spread of the
\emph{measured error} over tie-breakings of a zero-amplitude cut. It is a different quantity from the
$6.6\times10^{-4}$ by which ties move the \emph{bound} in the sweep of Sec.~\ref{sec:frontier}, where
the cut falls below the support floor and only two determinants are tied.) And the error at fraction $0.85$ has two independent determinations that differ by a
factor $2.3$, so only its order of magnitude is quoted. None of this touches the conclusion: the
two-level counterexample is analytic, and the counting argument of the next paragraph but one uses no
ranking, no protocol and no numerics at all. The sweep is corroboration of a conclusion that does not
depend on it.

\paragraph{Why the weight-only bound fails, in full.}
The reconstruction is Rayleigh--Ritz on $H_S=P_SHP_S$ (Sec.~\ref{sec:method}), not the restriction of
$A$ to a subset of its own eigenpairs: the Ritz energies $\tilde E_m$ are not eigenvalues of $H$, so
weight the subspace \emph{keeps} is nonetheless moved in $\omega$. The obvious sketch elides exactly
this. Having correctly noted $\int A=\lVert\phi\rVert^{2}$ and $\int A_S=\lVert P_S\phi\rVert^{2}$, it
bounds the $L_1$ discrepancy by twice the missing weight --- a step legitimate when $[P_S,H]=0$, where
$A_S$ is $A$ with the outside contributions deleted and
$\lVert A-A_S\rVert_1=\lVert\phi\rVert^{2}(1-w_S)$ identically, and controlling nothing otherwise. The
two-level example is the extreme: nothing is discarded there, and every unit of retained weight is
displaced by $g$. What the integrals do give, with no commutator hypothesis, is
$\lvert\int A-\int A_S\rvert\le\lVert A-A_S\rVert_{1}$, that is
$\lVert\phi\rVert^{2}(1-w_S)\le\lVert A-A_S\rVert_{1}$ --- a \emph{lower} bound, with $0$ violations across the $606$ subspaces of the pooled sweep. It survives intact as the
left-hand half of Theorem~\ref{thm:leakage}: the captured weight does say how wrong the reconstruction
must be, it never said how wrong it can be.

\paragraph{The normalization of the witness errors.} (Relative here and throughout this subsection means divided by $\lVert\phi\rVert^{2}$, the normalization
of Theorem~\ref{thm:leakage}; dividing instead by the in-window integral of the exact $A$, as the
repository's own relative $L_1$ does, moves them by about $1\%$ at $\eta=0.05\,t$ and by $2.6$ to
$3.3\%$ at $\eta=0.18\,t$, measured in Sec.~\ref{app:hyp}, and the two conventions are not interchanged
anywhere in this paper without saying so.)
